\chapter*{Bilan periodique - Fin du stage}
\addcontentsline{toc}{chapter}{Bilan periodique - Fin du stage}
\markboth{Bilan periodique - Fin du stage}{Bilan periodique - Fin du stage}

Ce bilan de fin de stage presente l'avancement reel realise sur le projet
MES-X Dependency Management au sein de Forvia Informatique Tunisie. Le travail
effectue a couvert l'ensemble de la chaine, depuis la collecte des artefacts
Azure DevOps jusqu'a la mise a disposition d'une vue de decision exploitable par
les equipes de release.

\section*{1. Objectifs de stage}

Les objectifs fixes au demarrage etaient les suivants:

\begin{itemize}
  \item centraliser la tracabilite entre work items, pull requests, commits,
  branches et tags;
  \item reduire le temps d'analyse pre-release en remplacant la navigation
  manuelle entre ecrans Azure DevOps;
  \item fournir une interface unique pour l'analyse unitaire, batch et par
  sprint;
  \item proposer une aide a la revue de commits via LLM, sans impacter la
  decision deterministe de readiness.
\end{itemize}

\section*{2. Travaux realises}

\textbf{Conception et architecture}
\begin{itemize}
  \item Definition d'une architecture modulaire: Sprint, Work Item, Pull
  Request, Commit, Decision, LLM.
  \item Formalisation du flux de bout en bout: selection $\rightarrow$
  extraction $\rightarrow$ enrichissement $\rightarrow$ aggregation
  $\rightarrow$ decision.
\end{itemize}

\textbf{Implementation backend (NestJS + TypeScript)}
\begin{itemize}
  \item Mise en place des endpoints REST et du streaming SSE pour le suivi de
  progression.
  \item Integration des APIs Azure DevOps (Work Item API et Git API).
  \item Reconstruction du graphe de dependances avec prevention des cycles,
  dedoublonnage et traitement batch.
  \item Enrichissement des commits (branches, tags, contexte repository) pour
  l'analyse de readiness.
  \item Ajout d'un module LLM assiste (declenchement a la demande, gestion des
  erreurs, maintien d'une validation humaine obligatoire).
  \item Structuration de l'observabilite via journalisation applicative
  (log/warn/error/debug).
\end{itemize}

\textbf{Implementation frontend (Vue 3 + Quasar + Vite)}
\begin{itemize}
  \item Ecran d'analyse unitaire d'un work item et mode batch multi-elements.
  \item Ecran d'exploration des sprints avec chargement progressif des work
  items.
  \item Vue de decision consolidee par repository avec etat de readiness.
  \item Integration de l'action de revue LLM directement depuis la vue de
  decision.
  \item Amelioration de l'experience utilisateur via retour temps reel durant
  les traitements longs.
\end{itemize}

\textbf{Qualite, integration et delivery}
\begin{itemize}
  \item Tests unitaires backend sur services et controleurs des modules
  critiques.
  \item Validation de scenario de bout en bout cote API.
  \item Industrialisation de la build/deploiement via pipeline CI/CD et image
  Docker.
\end{itemize}

\section*{3. Resultats obtenus}

Les resultats atteints en fin de stage sont:

\begin{itemize}
  \item une plateforme de tracabilite fonctionnelle et deployee;
  \item une vision unifiee des dependances cross-repositories a partir d'un
  point d'entree metier (work item ou sprint);
  \item une reduction des analyses manuelles repetitives grace a
  l'automatisation de l'aggregation;
  \item une meilleure fiabilite des decisions de promotion via une vue
  consolidee par repository.
\end{itemize}

\section*{4. Difficultes rencontrees et solutions apportees}

\begin{itemize}
  \item \textbf{Heterogeneite des donnees Azure DevOps:} mise en place d'une
  normalisation DTO pour stabiliser les traitements.
  \item \textbf{Complexite du graphe de dependances:} adoption d'une traversal
  avec dedoublonnage et controle des cycles.
  \item \textbf{Performance sur gros volumes:} regroupement des appels API,
  traitements batch et limitation des recalculs.
  \item \textbf{Lisibilite du suivi d'execution:} ajout de logs structures et
  d'un retour de progression cote UI.
\end{itemize}

\section*{5. Competences developpees}

Au terme du stage, les competences suivantes ont ete consolidees:

\begin{itemize}
  \item conception d'architecture logicielle modulaire et orientee services;
  \item developpement full-stack TypeScript (NestJS/Vue/Quasar);
  \item integration d'APIs enterprise (Azure DevOps) et gestion des contrats de
  donnees;
  \item mise en place de pratiques de qualite (tests, logs, validation
  incrementale);
  \item execution projet en mode Agile Scrum avec livraison iterative.
\end{itemize}

\section*{6. Suite proposee}

Pour la continuation du projet, les evolutions les plus pertinentes sont:

\begin{itemize}
  \item etendre les regles de decision a des cas release plus complexes
  (strategies multi-branches et multi-tags);
  \item renforcer la couverture de tests end-to-end automatisee cote frontend;
  \item ajouter un assistant de planification proposant des chemins de mise a
  niveau repository par repository a partir d'un work item.
\end{itemize}

\section*{7. Auto-evaluation de fin de stage}

Les objectifs principaux du stage sont atteints. La valeur apportee est
operationnelle (centralisation, acceleration d'analyse, support a la decision)
et technique (base logicielle modulaire et evolutive). Les axes d'amelioration
restent principalement lies a l'industrialisation avancee des tests et a
l'enrichissement progressif des regles metier de readiness.
